\documentclass[11pt,a4paper]{article}

\usepackage{fontspec}
\usepackage{polyglossia}
\setmainlanguage{hebrew}
\setotherlanguage{english}

% קומפילציה: xelatex federated_learning_carlton.tex
\newfontfamily\hebrewfont{David}[
  Path=C:/Windows/Fonts/,
  Extension=.ttf,
  UprightFont=david,
  BoldFont=davidbd
]
\setmainfont{David}[
  Path=C:/Windows/Fonts/,
  Extension=.ttf,
  UprightFont=david,
  BoldFont=davidbd,
  Script=Hebrew
]
\newfontfamily\englishfont{Latin Modern Roman}
\newfontfamily\hebrewfonttt{David}[
  Path=C:/Windows/Fonts/,
  Extension=.ttf,
  UprightFont=david,
  BoldFont=davidbd,
  Script=Hebrew
]

\usepackage{amsmath,amssymb,amsfonts}
\usepackage{mathtools}
\usepackage{booktabs}
\usepackage{tabularx}
\usepackage{enumitem}
\usepackage{geometry}
\usepackage{hyperref}
\usepackage{bidi}
\usepackage{float}
\usepackage{caption}
\usepackage{placeins}
\usepackage{xcolor}
\usepackage{listings}
\usepackage[framemethod=default]{mdframed}
\usepackage{tikz}
\usetikzlibrary{arrows.meta,positioning,fit,backgrounds,decorations.pathmorphing}

% תיבות ידידותיות ל-RTL (bidi) --- מבוססות mdframed (במקום tcolorbox)
% #1 = צבע (ברירת מחדל gray), #2 = כותרת
\newenvironment{infobox}[2][gray]{%
  \par\smallskip
  \begin{mdframed}[
    linecolor=#1!60!black,
    linewidth=0.7pt,
    backgroundcolor=#1!6,
    roundcorner=3pt,
    innertopmargin=7pt,
    innerbottommargin=7pt,
    innerleftmargin=8pt,
    innerrightmargin=8pt,
    skipabove=6pt,
    skipbelow=6pt
  ]%
  {\color{#1!70!black}\bfseries #2}\par\smallskip
}{%
  \end{mdframed}%
  \par\smallskip
}

\newenvironment{stepbox}[1]{\begin{infobox}[gray]{#1}}{\end{infobox}}

% תיבת קוד/פסאודו-קוד --- התוכן LTR כדי שלא יתהפך ב-RTL
\newenvironment{codebox}[1]{\begin{infobox}[gray]{#1}\begin{english}\setLTR}{\end{english}\end{infobox}}

\lstset{
  basicstyle=\ttfamily\small,
  breaklines=true,
  frame=none,
  columns=fullflexible,
  keepspaces=true,
  showstringspaces=false
}

\geometry{margin=2.5cm}
\hypersetup{
  colorlinks=true,
  linkcolor=blue!70!black,
  citecolor=blue!70!black,
  urlcolor=blue!70!black
}

\newcommand{\loss}{\ell}
\newcommand{\FedAvg}{\textsc{FedAvg}}
\newcommand{\FedProx}{\textsc{FedProx}}

\title{למידה מאוחדת (Federated Learning) ברשתות רדיו קוגניטיביות --- CARLTON}
\author{}
\date{}

\begin{document}
\maketitle
\tableofcontents
\newpage

%==============================================================================
\section{השיטה הקלאסית/הריכוזית (Centralized ML)}
%==============================================================================

במצב הרגיל, אנחנו מניחים שיש לנו בסיס נתונים אחד גדול ($D$) שיושב על שרת יחיד (Data Lake בענן).
המטרה שלנו היא למצוא את המשקולות $w$ של הרשת שיביאו למינימום את פונקציית ההפסד (השגיאה) הממוצעת על כל המידע:
\begin{equation}
  \min_{w} F(w) = \frac{1}{|D|} \sum_{(x,y) \in D} \loss(w; x, y)
  \label{eq:centralized}
\end{equation}

\subsection{למה זה רע לרשת הרדיו שלך (CARLTON)?}

\begin{description}[style=nextline, leftmargin=2.5cm, labelwidth=2.2cm]
  \item[\textbf{רוחב פס מטורף (Bandwidth)}]
    כדי לאמן מודל כזה, כל מנהל רשת בשטח צריך להזרים בכל מילי-שנייה את נתוני ה-SINR הגולמיים, תדרי השידור, ומטריצות הרעש לענן.
    ברשתות אלחוטיות זה אבסורד --- אנחנו מבזבזים את רוחב הפס היקר של הרשת רק כדי לשלוח נתונים על רוחב הפס!

  \item[\textbf{אבטחה ופרטיות (Privacy)}]
    נתוני הרדיו חושפים מיקומים מדויקים של משתמשים, דפוסי תנועה צבאיים או אסטרטגיות מסחריות של מפעילים (כמו סלקום ופרטנר).
    אף מפעיל לא יסכים להוציא את המידע הגולמי הזה מחוץ לשרת המקומי שלו.
\end{description}

%==============================================================================
\section{הגישה של Federated Learning: המהפכה המתמטית}
%==============================================================================

במקום לאסוף את המידע, אנחנו מבזרים את האופטימיזציה.
המידע הגולמי ($D_k$) נשאר ``נעול'' בתוך כל סוכן/מנהל רשת $k$.

הנוסחה הגלובלית החדשה שרוצים למזער נראית כך:
\begin{equation}
  \min_{w} F(w) = \sum_{k=1}^{K} \frac{n_k}{n} F_k(w)
  \label{eq:global-fl}
\end{equation}

כאשר פונקציית ההפסד המקומית של כל סוכן בנפרד היא:
\begin{equation}
  F_k(w) = \frac{1}{n_k} \sum_{(x,y) \in D_k} \loss(w; x, y)
  \label{eq:local-loss}
\end{equation}

\subsection{מה הנוסחה הזו אומרת בעצם?}

המתמטיקה מוכיחה שאנחנו מנסים להגיע לאותה תוצאה בדיוק כאילו כל המידע היה במקום אחד, אבל בדרך חכמה יותר:

\begin{itemize}
  \item השרת המרכזי (האגרגטור) \textbf{לא רואה} דגימות מידע $(x, y)$.
  \item כל סוכן $k$ מחשב את ה-Loss המקומי שלו ($F_k$) על המכשיר שלו, מעדכן את המשקולות שלו, ושולח לשרת רק את \textbf{השינוי במשקולות}.
  \item השרת מחבר את כל המשקולות בצורה משוקללת (לפי גודל המידע $\frac{n_k}{n}$) כדי ליצור את המודל הגלובלי המאוחד.
\end{itemize}

%==============================================================================
\section{למה זה קריטי למחקר שלך (Wireless \& CARLTON)}
%==============================================================================

הטקסט מסכם בצורה מדויקת שלוש סיבות שבגללן השילוב של CARLTON עם FL הוא פריצת דרך למחקר שלך:

\begin{enumerate}
  \item \textbf{חיסכון עצום בתקשורת (Bandwidth):}
    משקולות של רשת עצבית קטנה (כמו ה-DeepMellow שלך) שוקלות כמה עשרות או מאות קילובייטים (KB).
    לשלוח קובץ כזה פעם ב-$100$ צעדים לשרת זה ``חינמי'' לחלוטין מבחינת עומס על הרשת, בהשוואה לשליחה של גיגה-בייטים (GB) של מדידות רדיו וספקטרום רציפות.

  \item \textbf{עמידה ברגולציה ופרטיות:}
    המידע המקומי (ה-Replay Buffer של ה-SINR) לעולם לא עוזב את השרת של מנהל הרשת המקומית.
    הגנה מושלמת על פרטיות המשתמשים וסודיות הרשת.

  \item \textbf{ביזור עומסי חישוב (Scalability):}
    במקום שרת ענן אחד יקר שיצטרך לעבד מיליוני אפיזודות של RL מכל הארץ, החישוב הפיזי (הסימולציה והרצת הרשת על ה-SINR) מתחלק בין כל מנהלי הרשתות בשטח.
    השרת המרכזי צריך רק לעשות פעולת חיבור וממוצע פשוטה, שהיא מהירה וזולה מאוד.
\end{enumerate}

%==============================================================================
\section{המילון: מה אומרים הסימנים בפרויקט?}
%==============================================================================

\begin{description}[style=nextline, leftmargin=3cm, labelwidth=2.8cm]
  \item[$K$ (מספר הקליינטים)]
    כמות רשתות התקשורת החופפות בסימולציה שלך (למשל, $K=3$ עבור סלקום, פרטנר והרשת הצבאית).

  \item[$w$ (המשקולות הגלובליות)]
    ה``מוח'' המשותף והמעודכן ביותר של ה-DeepMellow שיושב בשרת המרכזי.

  \item[$w_k$ (המשקולות המקומיות)]
    ה``מוח'' של רשת $k$, אחרי שהיא התאמנה קצת לבד על נתוני ה-SINR שלה.

  \item[$n_k$ (כמות המידע)]
    כמות דגימות ה-SINR (או כמות האפיזודות) שרשת $k$ אספה בבאפר המקומי שלה.

  \item[$E$ (Local Epochs)]
    כמה צעדי אימון מנהל הרשת עושה לבד בבית על הבאפר שלו, לפני שהוא רץ לספר לחבר'ה (לשרת).

  \item[$T$ (סך כל הסיבובים)]
    כמה פעמים בסך הכל השרת והרשתות הולכים להחליף ביניהם משקולות לאורך כל הריצה.
\end{description}

%==============================================================================
\section{אנטומיה של סבב תקשורת (Federated Round)}
%==============================================================================

בכל סבב תקשורת $t$ (מתוך $T$ סבבים), מתרחש הריקוד הקבוע הבא:

\begin{enumerate}[label=\textbf{\arabic*.}]
  \item \textbf{בחירת משתתפים (Client Selection):}
    השרת בוחר אילו רשתות משתתפות בסבב הנוכחי (בסימולציה שלך, לרוב נבחר את כולן, כלומר קבוצה $S$).

  \item \textbf{שידור (Broadcast):}
    השרת שולח את המשקולות הנוכחיות שלו ($w_t$) לכל מנהלי הרשתות.
    כולם מתחילים את הסבב עם אותו מוח בדיוק.

  \item \textbf{חישוב מקומי (Local Computation):}
    כאן הקוד של \texttt{Coordinator.py} נכנס לפעולה.
    כל מנהל רשת מריץ את הלופ במשך $E$ פעמים (Epochs), אוסף נתונים, ומעדכן את המודל המקומי שלו באמצעות \texttt{DeepMellow.learn()}.
    התוצאה היא משקולות חדשות: $w_k^{t+1}$.

  \item \textbf{העלאה (Upload):}
    מנהלי הרשתות שולחים לשרת (\texttt{TrainBlock.py}) רק את המשקולות החדשות שלהם.

  \item \textbf{שילוב (Aggregation):}
    השרת משתמש באלגוריתם שנקרא \FedAvg{} (קיצור של Federated Averaging).
    הוא עושה ממוצע משוקלל של המשקולות לפי כמות המידע $n_k$ של כל רשת:
    \begin{equation}
      w_{t+1} = \sum_{k \in S} \frac{n_k}{n} w_k^{t+1}
      \label{eq:fedavg-round}
    \end{equation}
\end{enumerate}

%==============================================================================
\section{הדילמה: תקשורת מול חישוב (Client Drift)}
%==============================================================================

כדי שהמודל יתכנס, הוא צריך הרבה מידע.
יש לנו שתי דרכים להשיג את זה, וכל אחת מגיעה עם מחיר:

\begin{description}[style=nextline, leftmargin=3.2cm, labelwidth=3cm]
  \item[\textbf{אפשרות א' (תקשורת גבוהה)}]
    לעשות אימון מקומי קצר מאוד ($E=1$) ולשלוח מיד לשרת.
    המשמעות: הרשתות כל הזמן מסונכרנות, אבל אנחנו מציפים את ערוץ התקשורת בהעברות קבצים (משקולות של רשתות עצביות יכולות לשקול המון).

  \item[\textbf{אפשרות ב' (חישוב גבוה)}]
    להגיד למנהלי הרשתות להתאמן המון זמן לבד ($E$ גדול) ורק אז לשלוח לשרת.
    זה חוסך המון רוחב פס בתקשורת.
\end{description}

\subsection{הסכנה: Client Drift (סטיית קליינט)}

כאן יש מלכודת מתמטית.
אם נגדיר $E$ גדול מדי, כל מנהל רשת יתאמן יותר מדי זמן על בעיות ה-SINR הספציפיות שלו.

ברשת קוגניטיבית זה מסוכן: רשת א' תפתח ``תובנות קיצוניות'' שטובות רק לה, וכאשר השרת ינסה לעשות ממוצע בין התובנות הקיצוניות של רשת א' לתובנות הקיצוניות של רשת ב', הממוצע יצא מודל גרוע שלא מתאים לאף אחת מהן!
זה נקרא \textbf{Client Drift}.

%==============================================================================
\section{האתגר הגדול: Non-IID Data (מידע לא אחיד)}
%==============================================================================

בסטטיסטיקה, \textbf{IID} אומר שהמידע של כולם מתפלג באותו אופן.

אם הסביבה הייתה IID, כל רשת הייתה חווה בדיוק את אותן הפרעות ואותם משתמשים ראשיים (PUs).
במצב כזה, הממוצע של \FedAvg{} עובד כמו קסם ומתכנס למינימום המושלם.

אבל במציאות של רדיו קוגניטיבי, המידע הוא \textbf{Non-IID}.
רשת אחת יושבת ליד שדה תעופה וחווה הפרעות ממכ``ם בערוץ $1$, ורשת שנייה יושבת במרכז העיר וחווה עומס של טלפונים בערוץ $4$.
הנוסחאות המקומיות שלהן (פונקציות ההפסד $\loss$) מושכות לכיוונים שונים לחלוטין.

מכיוון שהמידע הוא Non-IID, לעיתים ממוצע פשוט (\FedAvg) לא יספיק, וזה בדיוק המקום שבו המחקר שלך יכול להתרחב בעתיד לאלגוריתמים מתקדמים יותר כמו \textbf{FedProx} (שמוסיף עונש מתמטי לסוכנים כדי שלא יתרחקו יותר מדי מהמודל הגלובלי).

עכשיו, כשהבנו את הדינמיקה הזו של סבבי התקשורת ואת הסכנה של Client Drift בסביבה עם הפרעות רדיו שונות (Non-IID), אנחנו מבינים בדיוק למה אנחנו צריכים את הפרמטרים האלו בקוד.

%==============================================================================
\section{האתגר המרכזי: נתונים IID מול Non-IID}
%==============================================================================

בעולם מושלם, הלמידה המאוחדת הייתה קלה מאוד.
בעולם האמיתי (וברדיו בפרט), הנתונים שכל מנהל רשת רואה הם שונים לחלוטין.

\begin{description}[style=nextline, leftmargin=2.8cm, labelwidth=2.5cm]
  \item[\textbf{IID}]
    (מידע אחיד ומושלם) --- כל מנהלי הרשתות רואים בדיוק את אותה התפלגות של הפרעות ומשתמשים.
    אם הנתונים היו IID, ממוצע פשוט של המשקולות היה עובד כמו קסם בכל פעם.

  \item[\textbf{Non-IID}]
    (מידע הטרוגני/מגוון) --- כל מנהל רשת חי ב``מציאות'' משלו.
    מנהל רשת א' חווה רק הפרעות (מקביל למודל שרואה רק תמונות של כלבים), בעוד מנהל רשת ב' רואה ערוצים פנויים לחלוטין (מודל שרואה רק תמונות של חתולים).
\end{description}

כדי להבין איך זה פוגש אותך בקוד, הנה פירוט סוגי ה-Non-IID המרכזיים שקיימים בסביבת הרדיו שלך (CARLTON):

\begin{table}[H]
  \centering
  \caption{סוגי Non-IID בסביבת הרדיו הקוגניטיבית}
  \label{tab:non-iid}
  \begin{tabular}{@{}>{\raggedright\arraybackslash}p{2.8cm}>{\raggedright\arraybackslash}p{4.2cm}>{\raggedright\arraybackslash}p{6.2cm}@{}}
    \toprule
    \textbf{סוג ה-Non-IID} & \textbf{המשמעות הכללית} & \textbf{מה זה אומר בפרויקט הרדיו (CARLTON)?} \\
    \midrule
    Label skew &
    הטיה בתיוגים &
    רשת אחת יושבת באזור שקט (ערוצים עם SINR גבוה), ורשת אחרת באזור פקוק. \\
    \addlinespace
    Feature skew &
    הטיה במאפיינים &
    ערוץ $3$ נראה שונה פיזיקלית (מבחינת תבנית הדעיכה) באזורים גיאוגרפיים שונים. \\
    \addlinespace
    Quantity skew &
    הבדל בכמות המידע &
    לרשת א' יש $15$ משתמשים (המון דגימות), ולרשת ב' יש רק $2$ משתמשים (מעט דגימות). \\
    \addlinespace
    Concept drift &
    הסביבה עצמה משתנה &
    בזמן האימון, רשתות אחרות מחליפות תדרים, ולכן מפת ההפרעות משתנה ברגע. \\
    \bottomrule
  \end{tabular}
\end{table}

%==============================================================================
\section{הפתרון: אלגוריתם FedAvg (Federated Averaging)}
%==============================================================================

כדי להתמודד עם הבעיות האלו, משתמשים באלגוריתם הקלאסי ביותר של גוגל: \FedAvg.

במקום שכל סוכן ישלח עדכון לשרת על כל דגימת רדיו בודדת (מה שיחנוק את רוחב הפס), כל סוכן מתאמן באופן מקומי למשך מספר אפיזודות (שנקראות כאן $E$), ורק בסוף מעלה את המשקולות הסופיות לשרת.

\subsection{למה ממוצע ``משוקלל'' (Weighted Averaging)?}

נניח שמנהל רשת א' אסף $10{,}000$ דגימות ב-Replay Buffer שלו, ומנהל רשת ב' אסף רק $100$.
כשהשרת עושה ממוצע, הוא לא יכול לתת להם משקל שווה של $50/50$.
הוא נותן משקל גדול יותר לרשת שלמדה מיותר נתונים.
הנוסחה המתמטית של השרת נראית כך:
\begin{equation}
  w_{t+1} = \sum_{k \in S_t} \frac{n_k}{\sum_{j \in S_t} n_j} w_k^{t+1}
  \label{eq:fedavg-weighted}
\end{equation}
כאשר $n_k$ הוא כמות הדגימות של סוכן $k$.

%==============================================================================
\section{פירוק האלגוריתם צעד אחר צעד}
\label{sec:algorithm-breakdown}
%==============================================================================

בוא נפרק את האלגוריתם צעד אחר צעד, בצורה הכי ברורה ופרקטית שיש.
נראה בדיוק מה קורה \textbf{בצד השרת}, מה קורה \textbf{בצד הסוכנים} (מנהלי הרשתות), ואיך הלופים האלו עובדים \textbf{במקביל}.

האלגוריתם מורכב משני חלקים: \textbf{אתחול} ו\textbf{לופ של סבבי תקשורת}.

\subsection{מבט על: מי עושה מה?}

\begin{infobox}[blue]{זרימת סבב FL אחד --- מי עושה מה?}
\textbf{1. שרת מרכזי} (\texttt{TrainBlock.py}): בחירת $S_t$ $\to$ שידור $w^t$ $\to$ המתנה $\to$ אגרגציה $w^{t+1}$

\vspace{6pt}
\hrule
\vspace{6pt}

\textbf{2. שידור למטה} --- כל הסוכנים מקבלים $w^t$ (משקולות, KB)

\vspace{6pt}

\textbf{3. אימון מקומי במקביל} (כל הסוכנים בו-זמנית, בלי לדבר ביניהם):
\begin{itemize}[nosep,leftmargin=1.2cm]
  \item \textbf{סוכן $k{=}1$} (\texttt{Coordinator.py}): $E$ epochs על $D_1$
  \item \textbf{סוכן $k{=}2$} (\texttt{Coordinator.py}): $E$ epochs על $D_2$
  \item \textbf{סוכן $k{=}K$} (\texttt{Coordinator.py}): $E$ epochs על $D_K$
\end{itemize}

\vspace{6pt}
\hrule
\vspace{6pt}

\textbf{4. העלאה למעלה} --- כל סוכן שולח $w_k$ + $n_k$ (משקולות בלבד)

\vspace{6pt}

\textbf{5. אגרגציה בשרת:}
$w^{t+1} = \sum_{k \in S_t} \frac{n_k}{n}\, w_k$
\end{infobox}

\begin{table}[H]
  \centering
  \caption{חלוקת אחריות בין שרת לסוכנים}
  \label{tab:server-vs-clients}
  \begin{tabular}{@{}>{\raggedright\arraybackslash}p{2.4cm}>{\raggedright\arraybackslash}p{5.8cm}>{\raggedright\arraybackslash}p{5.8cm}@{}}
    \toprule
    \textbf{שלב} & \textbf{שרת (\texttt{TrainBlock})} & \textbf{סוכנים (\texttt{Coordinator})} \\
    \midrule
    אתחול &
    יוצר $w^0$ אקראי &
    מחזיקים Replay Buffer מקומי $D_k$ \\
    \addlinespace
    סבב $t$ &
    בוחר $S_t$, משדר $w^t$, מחכה, ממצע &
    מורידים $w^t$, מתאמנים $E$ פעמים, מעלים $w_k$ \\
    \addlinespace
    סיום &
    שומר $w^T$ ב-\texttt{Train\_weights\_frl/} &
    ממשיכים לפעול עם המודל המעודכן \\
    \bottomrule
  \end{tabular}
\end{table}

%------------------------------------------------------------------------------
\subsection{שלב 0: האתחול (שרת)}
\label{subsec:step0-init}
%------------------------------------------------------------------------------

\begin{codebox}{פסאודו-קוד --- אתחול}
\begin{lstlisting}[language={}]
Server initializes w^0
\end{lstlisting}
\end{codebox}

השרת יוצר את רשת ה-DeepMellow הראשונית עם משקולות אקראיות לחלוטין (נסמן אותן כ-$w^0$).
ברגע זה, הרשת \textbf{לא יודעת כלום} על רדיו קוגניטיבי או על SINR.

\begin{description}[style=nextline, leftmargin=2.8cm, labelwidth=2.5cm]
  \item[\textbf{בקוד}]
    יצירת אובייקט \texttt{DeepMellow} חדש ב-\texttt{TrainBlock.py}.
    המשקולות נשמרות ב-\texttt{state\_dict} של PyTorch.

  \item[\textbf{מה לא קורה}]
    אף נתון SINR לא נשלח לשרת.
    כל סוכן כבר אוסף נתונים מקומית לתוך $D_k$ שלו (Replay Buffer), אבל השרת לא רואה אותם.
\end{description}

%------------------------------------------------------------------------------
\subsection{שלב 1: הלופ הגלובלי (שרת)}
\label{subsec:step1-global-loop}
%------------------------------------------------------------------------------

\begin{codebox}{פסאודו-קוד --- לופ גלובלי (שרת)}
\begin{lstlisting}[language={}]
for each round t = 0, 1, 2, ..., T-1:
    Select clients S_t subset of {1, ..., K}
    Broadcast w^t to all k in S_t
    Wait for uploads {w_k} from all k in S_t
    w^{t+1} <- weighted_average({w_k}, {n_k})
\end{lstlisting}
\end{codebox}

השרת מתחיל לנהל את סבבי התקשורת ($T$).
בכל סבב $t$, הוא בוחר קבוצה של קליינטים ($S_t$) מתוך כלל הרשתות ($K$).

בפרויקט CARLTON, מכיוון שיש מספר קטן של רשתות (למשל $K{=}3$ רשתות חופפות: סלקום, פרטנר, צבא),
השרת יבחר \textbf{תמיד את כל הרשתות} בכל סבב:
\[
  S_t = \{1, 2, \ldots, K\} \quad \text{לכל } t
\]

\textbf{חשוב:} בשלב זה השרת \emph{רק} משדר משקולות וממתין.
הוא \textbf{לא} מריץ אימון על SINR.

%------------------------------------------------------------------------------
\subsection{שלב 2: הלופ המקומי במקביל (הסוכנים)}
\label{subsec:step2-local-parallel}
%------------------------------------------------------------------------------

\begin{codebox}{פסאודו-קוד --- לופ מקומי (כל סוכן $k$, במקביל)}
\begin{lstlisting}[language={}]
for each client k in S_t in parallel:
    w_k <- w^t                          # download global weights
    for epoch e = 1 ... E:
        for batch B from D_k:
            w_k <- w_k - eta * grad(l(w_k; B))   # local SGD step
    upload w_k to server
\end{lstlisting}
\end{codebox}

כאן קורה הקסם: הפעולות הבאות קורות \textbf{במקביל} אצל כל מנהלי הרשתות בשטח (סלקום, פרטנר וכו'),
\textbf{מבלי שהם יודעים אחד על השני}.

\subsubsection{א. הורדת המודל הגלובלי}

\begin{codebox}{הורדה}
\begin{lstlisting}[language={}]
w_k <- w^t
\end{lstlisting}
\end{codebox}

כל סוכן $k$ מוריד מהשרת את המשקולות הגלובליות הנוכחיות ($w^t$) וטוען אותן לתוך הרשת שלו.
זוהי פקודת \texttt{load\_state\_dict()} ב-PyTorch.

\begin{itemize}
  \item לפני השלב: לכל הסוכנים אותו ``מוח'' בדיוק ($w^t$).
  \item הנתונים המקומיים ($D_k$) נשארים על המכשיר --- לא עוברים בערוץ.
\end{itemize}

\subsubsection{ב. לופ האפיזודות / האימון המקומי}

\begin{codebox}{אימון מקומי}
\begin{lstlisting}[language={}]
for epoch e = 1 ... E:
    for batch B from D_k:
        w_k <- w_k - eta * grad(l(w_k; B))
\end{lstlisting}
\end{codebox}

הסוכן מתנתק מהרשת החיצונית ומתחיל להתאמן ``בבית'' למשך $E$ פעמים (Epochs):

\begin{enumerate}[label=\alph*.]
  \item מריץ את רשת הרדיו באמצעות \texttt{Coordinator.py}.
  \item שולף \textbf{באצ'ים} ($B$) של נתוני SINR ותגמולים מתוך ה-Replay Buffer המקומי ($D_k$).
  \item מחשב את הגרדיאנט של פונקציית ההפסד:
    \[
      \nabla \loss(w_k; B)
    \]
  \item מעדכן את המשקולות המקומיות:
    \[
      w_k \leftarrow w_k - \eta \, \nabla \loss(w_k; B)
    \]
    ב-PyTorch זה הצעד הקלאסי של האופטימייזר: \texttt{optimizer.step()}.
\end{enumerate}

בסוף השלב הזה, לכל סוכן יש משקולות \textbf{קצת שונות} ($w_k$) שמותאמות בדיוק למה שהוא חווה באזור שלו.

\begin{infobox}[orange]{למה ``במקביל'' חשוב?}
סוכן 1 יכול להיות ליד שדה תעופה (הפרעות ממכ``ם), סוכן 2 במרכז עיר (עומס סלולרי), וסוכן 3 בצבא (דפוסי שידור שונים).
כולם מתאמנים \emph{בו-זמנית} על $D_k$ שונה --- אבל כולם מתחילים מאותו $w^t$.
\end{infobox}

\subsubsection{ג. העלאה לשרת}

כל סוכן מסיים את ה-$E$ אפיזודות שלו ושולח \textbf{רק} את המשקולות המעודכנות ($w_k$) חזרה לשרת המרכזי.
גודל ההעלאה: עשרות--מאות KB (משקולות DeepMellow), לא GB של מדידות SINR.

%------------------------------------------------------------------------------
\subsection{שלב 3: האגרגציה / שילוב התובנות (שרת)}
\label{subsec:step3-aggregation}
%------------------------------------------------------------------------------

\begin{codebox}{פסאודו-קוד --- אגרגציה (שרת)}
\begin{lstlisting}[language={}]
w^{t+1} <- sum over k in S_t of (n_k / sum_j n_j) * w_k
\end{lstlisting}
\end{codebox}

השרת חיכה שכולם יסיימו ויעלו את המשקולות שלהם.
כעת הוא מבצע את הממוצע המשוקלל (Weighted Average) כדי לייצר את המודל של הסבב הבא ($w^{t+1}$):
\begin{equation}
  w^{t+1} = \sum_{k \in S_t} \frac{n_k}{\sum_{j \in S_t} n_j} \, w_k
  \label{eq:fedavg-aggregation-step}
\end{equation}

\begin{itemize}
  \item רשת עם יותר דגימות ($n_k$ גדול) משפיעה יותר על המודל הגלובלי.
  \item אחרי האגרגציה, $w^{t+1}$ הופך ל-$w^t$ של הסבב הבא.
  \item הלופ חוזר עד $t = T-1$.
\end{itemize}

%------------------------------------------------------------------------------
\subsection{האלגוריתם המלא במבט אחד}
%------------------------------------------------------------------------------

\begin{codebox}{\FedAvg{} --- אלגוריתם מלא}
\begin{lstlisting}[language={}]
# --- Initialization (Server) ---
Server initializes w^0

# --- Communication Rounds ---
for each round t = 0, 1, 2, ..., T-1:

    # Server: client selection
    Select clients S_t subset of {1, ..., K}

    # Clients: parallel local training
    for each client k in S_t in parallel:
        w_k <- w^t
        for epoch e = 1 ... E:
            for batch B from D_k:
                w_k <- w_k - eta * grad(l(w_k; B))
        upload w_k to server

    # Server: aggregation
    w^{t+1} <- sum over k in S_t of (n_k / sum_j n_j) * w_k

# --- End ---
Server saves w^T
\end{lstlisting}
\end{codebox}

\subsection{ציר זמן של סבב יחיד}

\begin{figure}[H]
  \centering
  \begin{LTR}
  \begin{tikzpicture}[
    x=1.05cm, y=0.85cm,
    phase/.style={minimum height=0.75cm, align=center, font=\scriptsize, draw, rounded corners},
    srv/.style={phase, fill=blue!15},
    cli/.style={phase, fill=green!15},
    time/.style={font=\scriptsize, anchor=north}
  ]
    % time axis
    \draw[->, thick] (0,0) -- (12.5,0) node[right, font=\small] {זמן};

    % server row
    \node[left, font=\small\bfseries] at (-0.3, 2.2) {שרת};
    \node[srv, minimum width=1.8cm] (s1) at (1.2, 2.2) {בחירת $S_t$};
    \node[srv, minimum width=2.2cm] (s2) at (3.5, 2.2) {שידור $w^t$};
    \node[srv, minimum width=4.5cm, fill=blue!8] (s3) at (7.2, 2.2) {המתנה...};
    \node[srv, minimum width=2.5cm] (s4) at (11, 2.2) {אגרגציה $\to w^{t+1}$};

    % clients row
    \node[left, font=\small\bfseries] at (-0.3, 0.8) {סוכנים};
    \node[cli, minimum width=1.5cm, opacity=0.4] at (1.2, 0.8) {};
    \node[cli, minimum width=1.8cm] (c1) at (3.2, 0.8) {טעינת $w^t$};
    \node[cli, minimum width=5.5cm] (c2) at (7.5, 0.8) {$E$ epochs על $D_k$ (מקביל)};
    \node[cli, minimum width=2cm] (c3) at (11.5, 0.8) {העלאת $w_k$};

    % sync arrows
    \draw[dashed, gray] (s2.south) -- (c1.north);
    \draw[dashed, gray] (c3.north) -- (s4.south);

    \node[time] at (1.2, -0.35) {$t$};
    \node[time] at (11, -0.35) {$t{+}1$};
  \end{tikzpicture}
  \end{LTR}
  \caption{ציר זמן: השרת משדר, הסוכנים עובדים במקביל, ואז חוזרים לאגרגציה.}
  \label{fig:timeline-single-round}
\end{figure}

\subsection{מיפוי ישיר לקוד CARLTON}

\begin{table}[H]
  \centering
  \caption{מאיזה שלב באלגוריתם מגיע כל קטע קוד}
  \label{tab:code-mapping}
  \begin{tabular}{@{}>{\raggedright\arraybackslash}p{3cm}>{\raggedright\arraybackslash}p{4.2cm}>{\raggedright\arraybackslash}p{5.5cm}@{}}
    \toprule
    \textbf{שלב באלגוריתם} & \textbf{קובץ} & \textbf{פעולה בקוד} \\
    \midrule
    אתחול $w^0$ &
    \texttt{TrainBlock.py} &
    יצירת מודל DeepMellow ראשוני \\
    \addlinespace
    בחירת $S_t$ + שידור $w^t$ &
    \texttt{TrainBlock.py} &
    שליחת \texttt{state\_dict} לכל Worker \\
    \addlinespace
    $w_k \leftarrow w^t$ &
    \texttt{Worker.py} &
    \texttt{load\_state\_dict(global\_weights)} \\
    \addlinespace
    לופ $E$ epochs &
    \texttt{Coordinator.py} &
    איסוף SINR, תגמול, \texttt{DeepMellow.learn()} \\
    \addlinespace
    העלאת $w_k$ &
    \texttt{Worker.py} $\to$ \texttt{TrainBlock.py} &
    החזרת \texttt{state\_dict} מעודכן + $n_k$ \\
    \addlinespace
    אגרגציה $w^{t+1}$ &
    \texttt{TrainBlock.py} &
    \texttt{federated\_averaging()} \\
    \bottomrule
  \end{tabular}
\end{table}

%------------------------------------------------------------------------------
\subsection{סיכום צעד-אחר-צעד: שרת מול סוכנים}
\label{subsec:step-by-step-summary}
%------------------------------------------------------------------------------

האלגוריתם מורכב משני חלקים עיקריים: \textbf{אתחול חד-פעמי} ו\textbf{לופ של $T$ סבבי תקשורת}.
להלן כל שלב --- מי מבצע אותו, מה בדיוק קורה, ומה עובר בערוץ התקשורת.

\begin{stepbox}{שלב 0 --- אתחול (שרת)}
\textbf{מי:} שרת \quad|\quad \textbf{פעולה:} אתחול $w^0$\\[4pt]
\textbf{פירוט:} יצירת DeepMellow עם משקולות אקראיות. הרשת עדיין לא יודעת כלום על SINR.\\[4pt]
\textbf{תקשורת:} כלום
\end{stepbox}

\begin{stepbox}{שלב 1א --- בחירת משתתפים (שרת)}
\textbf{מי:} שרת \quad|\quad \textbf{פעולה:} בחירת $S_t$\\[4pt]
\textbf{פירוט:} בוחר אילו רשתות משתתפות. אצלך: תמיד $S_t=\{1,\ldots,K\}$ (סלקום, פרטנר, צבא).\\[4pt]
\textbf{תקשורת:} כלום
\end{stepbox}

\begin{stepbox}{שלב 1ב --- שידור (שרת $\to$ סוכנים)}
\textbf{מי:} שרת $\to$ כל הסוכנים \quad|\quad \textbf{פעולה:} שידור $w^t$\\[4pt]
\textbf{פירוט:} כל הסוכנים מקבלים אותו ``מוח'' גלובלי לתחילת הסבב.\\[4pt]
\textbf{תקשורת:} משקולות $w^t$ (KB)
\end{stepbox}

\begin{stepbox}{שלב 2א --- הורדת מודל (סוכן $k$)}
\textbf{מי:} סוכן $k$ \quad|\quad \textbf{פעולה:} $w_k \leftarrow w^t$\\[4pt]
\textbf{פירוט:} \texttt{load\_state\_dict()} --- טעינת המשקולות ל-DeepMellow המקומי.\\[4pt]
\textbf{תקשורת:} ---
\end{stepbox}

\begin{stepbox}{שלב 2ב --- אימון מקומי (סוכן $k$)}
\textbf{מי:} סוכן $k$ \quad|\quad \textbf{פעולה:} אימון $E$ epochs\\[4pt]
\textbf{פירוט:} לולאה על באצ'ים $B \subset D_k$: הרצת \texttt{Coordinator.py}, חישוב $\nabla\loss$, \texttt{optimizer.step()}.\\[4pt]
\textbf{תקשורת:} --- \quad (הסוכנים עובדים \textbf{במקביל}, בלי לדבר ביניהם)
\end{stepbox}

\begin{stepbox}{שלב 2ג --- העלאה (סוכן $k$ $\to$ שרת)}
\textbf{מי:} סוכן $k$ $\to$ שרת \quad|\quad \textbf{פעולה:} העלאת $w_k$\\[4pt]
\textbf{פירוט:} שליחת המשקולות \textbf{בלבד} אחרי $E$ אפיזודות. הנתונים הגולמיים נשארים מקומית.\\[4pt]
\textbf{תקשורת:} משקולות $w_k$ + $n_k$
\end{stepbox}

\begin{stepbox}{שלב 3 --- אגרגציה (שרת)}
\textbf{מי:} שרת \quad|\quad \textbf{פעולה:} אגרגציה\\[4pt]
\textbf{פירוט:} ממתין שכל $k \in S_t$ יסיים, ואז:
$w^{t+1} = \sum_k \frac{n_k}{n} w_k$\\[4pt]
\textbf{תקשורת:} כלום
\end{stepbox}

\begin{infobox}[green]{הלופ במקביל --- מה זה אומר בפועל?}
בשלב 2, \textbf{כל} מנהלי הרשתות (סלקום, פרטנר, צבא) מריצים את \texttt{Coordinator.py} \emph{בו-זמנית} על המחשבים שלהם.
הם לא מתקשרים ביניהם --- רק עם השרת (הורדה בתחילת הסבב, העלאה בסופו).
השרת בינתיים ``יושב ומחכה'' עד שכולם מסיימים את ה-$E$ אפיזודות שלהם.
\end{infobox}

\begin{infobox}[blue]{דוגמה מספרית לסבב אחד ($K=3$, $E=5$)}
\begin{enumerate}[label=\arabic*., nosep]
  \item \textbf{שרת} שולח $w^{42}$ לשלוש הרשתות.
  \item \textbf{סלקום} מתאמן 5 אפיזודות על $D_1$ (עיר צפופה, $n_1{=}8000$) $\Rightarrow$ מקבל $w_1$.
  \item \textbf{פרטנר} מתאמן 5 אפיזודות על $D_2$ (כפר שקט, $n_2{=}2000$) $\Rightarrow$ מקבל $w_2$.
  \item \textbf{צבא} מתאמן 5 אפיזודות על $D_3$ (שדה תעופה, $n_3{=}5000$) $\Rightarrow$ מקבל $w_3$.
  \item \textbf{שרת} מחשב: $w^{43} = \frac{8000}{15000}w_1 + \frac{2000}{15000}w_2 + \frac{5000}{15000}w_3$.
\end{enumerate}
\end{infobox}

%==============================================================================
\section{הדילמה הגדולה: תקשורת לעומת סטיית קליינט}
%==============================================================================

כאן אתה כמפתח צריך לקבל החלטה שמשפיעה על הביצועים של הפרויקט:

\begin{itemize}
  \item \textbf{$E$ גדול} (הרבה אימון מקומי):
    חוסך המון תקשורת (פחות סבבים), אבל מסתכן בכך שכל סוכן יפתח ``תובנות שגויות'' שמתאימות רק לאזור הגיאוגרפי שלו, והממוצע ייכשל (זה נקרא Client Drift).

  \item \textbf{$E$ קטן} (למשל, סבב אחרי כל אפיזודה בודדת):
    מונע סטיות ושומר על כולם מסונכרנים, אבל דורש המון העברות של משקולות ברשת הפיזית.
\end{itemize}

%==============================================================================
\section{\FedProx{} --- אופטימיזציה מאוחדת עם עונש קרבה}
\label{sec:fedprox}
%==============================================================================

\begin{quote}
  \textbf{Li et al., 2020} --- \textit{Federated Optimization in Heterogeneous Networks}.\\
  מטפל בהטרוגניות (Non-IID) ובסטיית קליינט (Client Drift).
\end{quote}

%------------------------------------------------------------------------------
\subsection{מוטיבציה: למה \FedAvg{} נכשל ב-Non-IID ואיך \FedProx{} מציל את המצב?}
\label{subsec:fedprox-motivation}
%------------------------------------------------------------------------------

כפי שראינו, ב-\FedAvg, אם לרשת א' ולרשת ב' יש בעיות רדיו שונות לחלוטין (Non-IID), האימון המקומי שלהן ימשוך את המשקולות שלהן ($w_k$) לכיוונים הפוכים ומרוחקים מאוד מהמודל הגלובלי ($w^t$).
כשהשרת יעשה ממוצע, התוצאה תהיה מודל גרוע שלא מתאים לאף אחת מהן.

\FedProx{} פותר את זה על ידי הוספת ``רסן'' או ``קנס מתמטי'' בזמן האימון המקומי בסוכנים.

\begin{description}[style=nextline, leftmargin=2.8cm, labelwidth=2.5cm]
  \item[\FedAvg]
    כל סוכן ממזער רק את $F_k(w)$ --- חופשי לנדוד לאן שירצה בתוך נוף ההפסד המקומי שלו.

  \item[\FedProx]
    מוסיף איבר קרבה שמעניש על מרחק מ-$w^t$ --- הסוכן לומד מהמידע המקומי, אבל לא ``בורח'' רחוק מדי מהקו הגלובלי.
\end{description}

\textbf{דוגמה מ-CARLTON:}
רשת א' ליד שדה תעופה (הפרעות ממכ``ם בערוץ 1) ורשת ב' במרכז עיר (עומס סלולרי בערוץ 4) יפתחו גרדיאנטים שמושכים את $w_k$ לכיוונים מנוגדים.
ב-\FedAvg, הממוצע של שני הכיוונים האלה יוצר מודל ``אמצע'' גרוע.
ב-\FedProx, כל סוכן נשאר קרוב יותר ל-$w^t$, והממוצע יציב ושימושי יותר.

%------------------------------------------------------------------------------
\subsection{פונקציית המטרה המקומית}
\label{subsec:fedprox-objective}
%------------------------------------------------------------------------------

במקום שסוכן $k$ ינסה רק למזער את פונקציית ההפסד הרגילה שלו ($F_k(w)$), ב-\FedProx{} הוא מנסה לפתור את המשוואה הבאה:
\begin{equation}
  \min_{w} \; F_k(w) + \frac{\mu}{2} \, \lVert w - w^t \rVert^2
  \label{eq:fedprox-objective}
\end{equation}

בוא נבין את שני החלקים של הנוסחה:

\begin{description}[style=nextline, leftmargin=3.2cm, labelwidth=3cm]
  \item[$F_k(w)$]
    פונקציית ההפסד הרגילה של הסוכן.
    במקרה שלך: טעות ה-Q-learning של ה-DeepMellow על נתוני ה-SINR המקומיים ב-$D_k$.

  \item[$\dfrac{\mu}{2} \lVert w - w^t \rVert^2$]
    \textbf{האיבר הפרוקסימלי (Proximal Term).}
    זהו המרחק הגיאומטרי (מרחק אוקלידי בריבוע) בין המשקולות הנוכחיות של הסוכן ($w$) לבין המשקולות הגלובליות שהשרת שלח לו בתחילת הסבב ($w^t$).
\end{description}

\subsubsection{מה עושה הפרמטר $\mu$ (מיו)?}

$\mu$ הוא כפתור ה``ווליום'' של הקנס:

\begin{table}[H]
  \centering
  \caption{השפעת הפרמטר $\mu$ ב-\FedProx}
  \label{tab:mu-effect}
  \begin{tabular}{@{}>{\raggedright\arraybackslash}p{2.2cm}>{\raggedright\arraybackslash}p{11.5cm}@{}}
    \toprule
    \textbf{ערך $\mu$} & \textbf{משמעות} \\
    \midrule
    $\mu = 0$ &
    האיבר הימני נמחק, והאלגוריתם חוזר להיות \FedAvg{} רגיל. הסוכן יכול לנדוד לאן שהוא רוצה. \\
    \addlinespace
    $\mu$ ענק &
    הסוכן יקבל ``עונש'' עצום אם הוא ישנה אפילו קצת את המשקולות. המודל כמעט ולא יזוז מהמקום (שמרני מדי). \\
    \addlinespace
    $\mu$ מכוון נכון &
    (למשל $0.1$ או $0.01$) --- נוצר איזון מושלם: הסוכן מצד אחד לומד מה-SINR המקומי, ומצד שני נשאר קרוב מספיק לקו המחשבה הגלובלי של השרת. \\
    \bottomrule
  \end{tabular}
\end{table}

\textbf{אינטואיציה מתמטית:} האיבר $\frac{\mu}{2}\lVert w - w^t \rVert^2$ מוסיף ``עלות'' לכל סטייה מ-$w^t$.
בעדכון גרדיאנט, זה שקול לכוח משיכה חזרה לעוגן הגלובלי:
\begin{equation}
  w_k \leftarrow w_k - \eta \left( \nabla F_k(w_k) + \mu \, (w_k - w^t) \right)
  \label{eq:fedprox-update}
\end{equation}

%------------------------------------------------------------------------------
\subsection{אינטואיציה: מודל ``הגומייה''}
\label{subsec:fedprox-rubber-band}
%------------------------------------------------------------------------------

חשוב על $w^t$ כ\textbf{עוגן} (anchor).
כל קליינט חוקר את נוף ההפסד המקומי שלו, אבל ``נמשך'' חזרה לעוגן.
זה מפחית עדכונים סותרים כשהגרדיאנטים מצביעים לכיוונים שונים (Non-IID).

\begin{figure}[H]
  \centering
  \begin{LTR}
  \begin{tikzpicture}[
    scale=1.05,
    anchorpoint/.style={circle, fill=blue!70, inner sep=2pt},
    endpoint/.style={circle, fill=red!70, inner sep=2pt},
    drift/.style={-{Stealth}, thick, dashed, red!70},
    prox/.style={-{Stealth}, thick, green!50!black},
    band/.style={green!50!black, thick, decorate, decoration={coil, aspect=0.3, segment length=3pt, amplitude=2pt}}
  ]
    % loss landscape contour (schematic)
    \draw[gray!40, thick] (0,0) ellipse (2.8cm and 1.6cm);
    \draw[gray!25] (-0.5,0.3) ellipse (2.2cm and 1.2cm);
    \node[font=\scriptsize, gray] at (0,-2.1) {נוף ההפסד המקומי $F_k(w)$};

    % anchor w^t
    \node[anchorpoint, label={above:$w^t$ (עוגן)}] (wt) at (0,0) {};
    \draw[blue!50, dashed] (wt) circle (1.1cm);
    \node[font=\scriptsize, blue!60] at (1.55,0.55) {אזור מותר ($\mu$)};

    % FedAvg drift - wanders far
    \node[endpoint, label={below right:\textcolor{red!70}{סטיית \FedAvg}}] (wfar) at (2.4,-0.9) {};
    \draw[drift] (wt) -- (wfar);

    % FedProx - stays near anchor
    \node[endpoint, label={above right:\textcolor{green!50!black}{\FedProx}}] (wnear) at (0.75,0.55) {};
    \draw[prox] (wt) -- (wnear);
    \draw[band] (wt) -- (wnear);

    % legend
    \node[font=\scriptsize, align=left, anchor=north west] at (-3.2,-2.5) {
      $\mu$ מגביל עד כמה האימון המקומי\\
      יכול לנדוד בכל סבב
    };
  \end{tikzpicture}
  \end{LTR}
  \caption{מודל הגומייה: \FedAvg{} עלול לנדוד רחוק מ-$w^t$; \FedProx{} מגביל את הסטייה.}
  \label{fig:fedprox-rubber-band}
\end{figure}

\begin{itemize}
  \item \textbf{ללא עונש} (\FedAvg): הסוכן ``רץ'' לכיוון המינימום המקומי שלו --- עלול להגיע רחוק מ-$w^t$.
  \item \textbf{עם עונש} (\FedProx): הגומייה מושכת חזרה --- הסוכן נשאר באזור סביב העוגן הגלובלי.
  \item \textbf{תוצאה:} הממוצע בשרת מקבל משקולות פחות סותרות, והאימון יציב יותר תחת Non-IID.
\end{itemize}

%------------------------------------------------------------------------------
\subsection{מתי נרצה להשתמש ב-\FedProx{} ברשת הרדיו?}
\label{subsec:fedprox-when}
%------------------------------------------------------------------------------

\begin{enumerate}[label=\textbf{\arabic*.}]
  \item \textbf{פרופילי הפרעה שונים (Heterogeneity):}
    אם רשת א' נמצאת באזור עירוני צפוף (המון הפרעות) ורשת ב' באזור כפרי שקט.
    ה-NON-IID כאן הוא לא רק סטטיסטי --- הוא פיזי: מפות SINR שונות לגמרי.

  \item \textbf{חוסר יציבות באימון:}
    אם במהלך הריצות של \FedAvg{} אתה רואה שגרף התגמול (Reward) או ה-Loss קופץ בפראות (מתנדנד) או מתרחק לחלוטין (לא מתכנס).
    הוספת $\mu$ תרגיע את הגרף ותייצב את הלמידה.

  \item \textbf{$E$ גדול:}
    כשאתה רוצה לחסוך תקשורת (אימון מקומי ארוך) אבל חושש מסטיית קליינט --- $\mu$ מאפשר $E$ גבוה יותר בבטחה.

  \item \textbf{השתתפות חלקית ונפילות רשת (Partial Participation):}
    בעולם הרדיו האמיתי, מכשירים או מנהלי רשתות יכולים להתנתק באמצע סבב (בגלל בעיות קליטה, סוללה או עומס).
    \FedProx{} מוכח מתמטית כאלגוריתם ששומר על יציבות המודל הגלובלי גם כאשר רק חלק מהסוכנים (שבר $C$) מצליחים להחזיר את המשקולות שלהם בסוף הסבב:
    \[
      |S_t| = \lfloor C \cdot K \rfloor, \quad 0 < C \leq 1
    \]
\end{enumerate}

%------------------------------------------------------------------------------
\subsection{פירוק האלגוריתם צעד אחר צעד (\FedProx{})}
\label{subsec:fedprox-algorithm-steps}
%------------------------------------------------------------------------------

\FedProx{} זהה ל-\FedAvg{} בכל השלבים --- \textbf{חוץ משלב 2ב} (האימון המקומי).
השרת, ההורדה, ההעלאה והאגרגציה נשארים בדיוק אותו דבר.

\begin{table}[H]
  \centering
  \caption{השוואת שלבים: \FedAvg{} מול \FedProx}
  \label{tab:fedavg-vs-fedprox-steps}
  \begin{tabular}{@{}c>{\raggedright\arraybackslash}p{5.5cm}cc@{}}
    \toprule
    \textbf{שלב} & \textbf{פעולה} & \FedAvg & \FedProx \\
    \midrule
    0 & אתחול $w^0$ (שרת) & $\checkmark$ & $\checkmark$ \\
    1א & בחירת $S_t$ (שרת) & $\checkmark$ & $\checkmark$ \\
    1ב & שידור $w^t$ (שרת) & $\checkmark$ & $\checkmark$ \\
    2א & $w_k \leftarrow w^t$ (סוכן) & $\checkmark$ & $\checkmark$ \\
    2ב & אימון $E$ epochs (סוכן) & $\nabla F_k$ בלבד & $\nabla F_k + \mu(w_k{-}w^t)$ \\
    2ג & העלאת $w_k$ (סוכן) & $\checkmark$ & $\checkmark$ \\
    3 & אגרגציה $w^{t+1}$ (שרת) & $\checkmark$ & $\checkmark$ \\
    \bottomrule
  \end{tabular}
\end{table}

\begin{codebox}{\FedProx{} --- אלגוריתם מלא (ההבדל מסומן ב-*)}
\begin{lstlisting}[language={}]
# --- Initialization (Server) ---
Server initializes w^0

# --- Communication Rounds ---
for each round t = 0, 1, 2, ..., T-1:

    # Server: client selection + broadcast
    Select clients S_t subset of {1, ..., K}
    Broadcast w^t to all k in S_t

    # Clients: parallel local training
    for each client k in S_t in parallel:
        w_k <- w^t
        for epoch e = 1 ... E:
            for batch B from D_k:
                # * FedProx only: proximal term pulls w_k toward w^t
                w_k <- w_k - eta * (grad(F_k(w_k; B)) + mu * (w_k - w^t))
        upload w_k to server

    # Server: aggregation (identical to FedAvg)
    w^{t+1} <- sum over k in S_t of (n_k / sum_j n_j) * w_k

# --- End ---
Server saves w^T
\end{lstlisting}
\end{codebox}

\begin{infobox}[orange]{מה משתנה בקוד CARLTON?}
רק שורת העדכון ב-\texttt{DeepMellow.learn()} / האופטימייזר:
במקום \texttt{loss.backward()} על $F_k$ בלבד, מוסיפים את האיבר $\frac{\mu}{2}\|w - w^t\|^2$ ל-loss לפני ה-backward.
האגרגציה ב-\texttt{TrainBlock.py} נשארת \texttt{federated\_averaging()} ללא שינוי.
\end{infobox}

%------------------------------------------------------------------------------
\subsection{השתתפות חלקית --- פירוט נוסף}
\label{subsec:fedprox-partial}
%------------------------------------------------------------------------------

\begin{description}[style=nextline, leftmargin=3cm, labelwidth=2.8cm]
  \item[\textbf{בעיה ב-\FedAvg}]
    כשמשתתף רק חלק מהרשתות, עדכונים ממקורות שונים בכל סבב --- האימון עלול להיות לא יציב.

  \item[\textbf{יתרון \FedProx}]
    העונש הקרבה מייצב את האימון גם כשההשתתפות מדורבנת --- מצב נפוץ ברשתות ניידות ואלחוטיות.

  \item[\textbf{ב-CARLTON}]
    מנהל רשת אחד עלול להיות לא זמין (תחזוקה, קליטה חלשה).
    \FedProx{} מאפשר להמשיך אימון עם $S_t \subset \{1,\ldots,K\}$ בלי שהמודל הגלובלי ``יתפרק''.
\end{description}

\begin{infobox}[blue]{השוואה קצרה: \FedAvg{} מול \FedProx{}}
\begin{center}
\begin{tabular}{@{}lcc@{}}
  \toprule
  & \FedAvg & \FedProx \\
  \midrule
  אגרגציה בשרת & ממוצע משוקלל & ממוצע משוקלל (זהה) \\
  אימון מקומי & $\min F_k(w)$ & $\min F_k(w) + \frac{\mu}{2}\lVert w-w^t\rVert^2$ \\
  פרמטר נוסף & --- & $\mu \geq 0$ \\
  Non-IID / Drift & רגיש & עמיד יותר \\
  השתתפות חלקית & פחות יציב & יציב יותר \\
  \bottomrule
\end{tabular}
\end{center}
\end{infobox}

%==============================================================================
\section{החיבור לקוד שלך}
%==============================================================================

בסעיף ההטמעה הספציפית, מוזכרת ההטמעה ב-\texttt{BuildingBlocks/TrainBlock.py}.
הפונקציה \texttt{federated\_averaging()} שתיצור ב-PyTorch תיקח את ה-\texttt{state\_dict} (מילון המשקולות של PyTorch) מכל מנהלי הרשתות, תבצע עליהם את הממוצע המשוקלל הזה, ולאחר $1000$ סבבים תשמור את התוצאה בתיקייה \texttt{Train\_weights\_frl/}.

מכיוון שאצלך כל אפיזודה יכולה להסתיים בזמן קצת שונה או עם כמות צעדים שונה, כמות הנתונים ($n_k$) של כל סוכן עשויה להשתנות.

\end{document}
